% !TeX program = lualatex
% ================================================================
% Urdu Parallel Text Version of PieChartsDrawing
%
% Generated by the server using the following settings:
%
%   language            Urdu (urd)
%   text direction      right to left
%   font                Noto Nastaliq Urdu
%   fallback font       Noto Serif
%   built from          latex/worksheets/charts/pieChartsDrawing.tex
%   vocabulary key      latex/worksheets/charts/pieChartsDrawing/L2/urd/pieChartsDrawing_urduKey.tex
%   tier 3 vocabulary   green   RGB 0 128 0
%   English text        black   RGB 0 0 0
%   Urdu text           purple  RGB 112 48 160
%   layout macros       version 3
%
% Please don't edit any information in this comment block - the server
% compares it against the current settings and rebuilds this file when
% they differ, so changes here would be overwritten. However, feel free
% to make updates to the LaTeX below if you want to edit it.
% ================================================================

\documentclass[a4paper,12pt]{article}
% ================================================================
% Copied from the English sheet this was translated from, so that
% anything it sets up for itself still works here
% ================================================================
\usepackage[a4paper, margin=1.8cm]{geometry}
\usepackage{tikz}
\usepackage{amsmath}
\usepackage{amssymb}
\usepackage{array}
\usepackage{mdframed}
\usepackage{enumitem}
\usepackage[T1]{fontenc}
\usepackage{textcomp}
\usetikzlibrary{calc}

% ---- Colours ----
\definecolor{slice1}{RGB}{70,130,180}
\definecolor{slice2}{RGB}{240,140,40}
\definecolor{slice3}{RGB}{70,170,70}
\definecolor{slice4}{RGB}{210,70,70}
\definecolor{slice5}{RGB}{150,90,190}
\definecolor{lightgrey}{RGB}{245,245,245}
\definecolor{boxblue}{RGB}{220,235,250}
\definecolor{boxgreen}{RGB}{220,245,220}
\definecolor{boxorange}{RGB}{255,240,220}

% -------------------------------------------------------
% \emptycircle{radius_cm} -- just a blank circle with centre dot and 12-o'clock line
% -------------------------------------------------------
\newcommand{\emptycircle}[1]{%
  \begin{tikzpicture}
    \pgfmathsetmacro{\r}{#1}%
    \draw[black, line width=1.8pt] (0,0) circle (\r cm);
    \fill[black] (0,0) circle (3pt);
    \draw[black, line width=1.5pt] (0,0) -- (90:\r cm);
    \draw[black, line width=1.2pt] (0,\r cm) -- (0,{\r cm+0.25cm});
  \end{tikzpicture}%
}

% -------------------------------------------------------
% \partialchart{radius}{Label/Angle/Colour, ...} -- draws some sectors pre-filled
%   remaining arc is left blank for the student to complete
%   sectors listed are drawn; the rest is empty
% -------------------------------------------------------
\newcommand{\partialchart}[2]{%
  \begin{tikzpicture}
    \pgfmathsetmacro{\r}{#1}%
    \pgfmathsetmacro{\startA}{90}%
    \foreach \lbl/\ang/\col in {#2} {%
      \pgfmathsetmacro{\endA}{\startA - \ang}%
      \fill[\col!25] (0,0) -- (\startA:\r cm) arc(\startA:\endA:\r cm) -- cycle;%
      \draw[black, line width=1.8pt] (0,0) -- (\startA:\r cm);%
      \draw[black, line width=1.8pt] (0,0) -- (\endA:\r cm);%
      \pgfmathsetmacro{\midA}{(\startA + \endA)/2}%
      \pgfmathsetmacro{\lR}{\r * 0.65}%
      \node[font=\small\bfseries, align=center] at (\midA:\lR cm) {\lbl};%
      \xdef\startA{\endA}%
    }%
    \draw[black, line width=1.8pt] (0,0) circle (\r cm);%
    \fill[black] (0,0) circle (3pt);%
    \draw[black, line width=1.2pt] (0,\r cm) -- (0,{\r cm+0.25cm});%
  \end{tikzpicture}%
}

\newcommand{\ansbox}{\underline{\hspace{1.6cm}}}

\pagestyle{empty}

% Characters this language's font has not got are borrowed from another,
% rather than being silently dropped - see docs/translated-sheet-latex.md
\directlua{%
  if luaotfload and luaotfload.add_fallback then
    luaotfload.add_fallback("eallatinfallback",
      { "Noto Serif:mode=harf;" })
  else
    tex.error("this luaotfload is too old to register a fallback font")
  end
}
\usepackage[english,bidi=basic]{babel}
\babelprovide[import]{urdu}
\babelfont[urdu]{rm}[Renderer=Harfbuzz, BoldFont={Noto Nastaliq Urdu}, ItalicFont={Noto Nastaliq Urdu}, BoldItalicFont={Noto Nastaliq Urdu}, RawFeature={fallback=eallatinfallback}]{Noto Nastaliq Urdu}
\babelfont[urdu]{sf}[Renderer=Harfbuzz, BoldFont={Noto Nastaliq Urdu}, ItalicFont={Noto Nastaliq Urdu}, BoldItalicFont={Noto Nastaliq Urdu}, RawFeature={fallback=eallatinfallback}]{Noto Nastaliq Urdu}
\newcommand{\ealtext}[1]{\foreignlanguage{urdu}{#1}}
\newcommand{\ealtextblock}[1]{{\raggedleft\foreignlanguage{urdu}{#1}\par}}
% ================================================================
% EAL layout helpers
% ================================================================
\usepackage{xcolor}

\definecolor{ealkeycolour}{RGB}{0,128,0}
\definecolor{ealenglish}{RGB}{0,0,0}
\definecolor{ealtwocolour}{RGB}{112,48,160}

\newsavebox{\ealboxen}
\newsavebox{\ealboxtr}
\newlength{\ealglosswd}
\newlength{\ealglossraise}
\setlength{\ealglossraise}{1.15em}

% a tier 3 word, in English and in translation alike
\newcommand{\ealkey}[1]{\textcolor{ealkeycolour}{#1}}
\newcommand{\ealkeytr}[1]{\textcolor{ealkeycolour}{\ealtext{#1}}}

% One translated word above one English word. Built from kernel
% primitives, never a tabular, which would break wherever the sheet
% loads array, booktabs or tabularx. The box is as wide as the wider of
% the two, so a long translation pushes its neighbours aside instead of
% overlapping them, and it claims its full height so a table row or a
% TikZ node makes room for it.
\newcommand{\ealgloss}[2]{%
  \sbox{\ealboxen}{\ealkey{#1}}%
  \sbox{\ealboxtr}{\small\textcolor{ealtwocolour}{\ealtext{#2}}}%
  \setlength{\ealglosswd}{\wd\ealboxen}%
  \ifdim\wd\ealboxtr>\ealglosswd\setlength{\ealglosswd}{\wd\ealboxtr}\fi
  \leavevmode
  \makebox[\ealglosswd][c]{%
    \makebox[0pt][c]{\usebox{\ealboxen}}%
    \makebox[0pt][c]{\raisebox{\ealglossraise}{\usebox{\ealboxtr}}}%
  }%
}

% opens up line spacing around a block of glosses, so the raised
% translations sit clear of the line above
\newenvironment{ealglossed}{\par\linespread{2.1}\selectfont\ignorespaces}{\par}

% a whole translated sentence above its English counterpart. block level,
% so it does not belong inside a tabular cell
\newcommand{\ealpara}[2]{%
  \par\addvspace{0.5\baselineskip}%
  {\color{ealtwocolour}\ealtextblock{#1}}%
  \penalty200
  {\color{ealenglish}#2\par}%
  \addvspace{0.5\baselineskip}%
}

\begin{document}

% ===========================
%  TITLE
% ===========================
\begin{center}
\ealpara{\ealkeytr{پائی چارٹ} \ealkeytr{زاویہ نما} کے ساتھ کھینچنا}{Drawing \ealkey{Pie Charts} with a \ealkey{Protractor}}

\ealpara{نام: \underline{\hspace{6cm}} \quad تاریخ: \underline{\hspace{2.8cm}}}{Name:\;\underline{\hspace{6cm}} \quad Date:\;\underline{\hspace{2.8cm}}}
\end{center}
\vspace{2pt}


\begin{mdframed}[backgroundcolor=boxgreen, linecolor=slice3,
                 linewidth=1.5pt, roundcorner=6pt, innertopmargin=6pt]
\ealpara{کلیدی فارمولا: \ealkeytr{قطاع} \ealkeytr{زاویہ} = \dfrac{مقدار}{کل} \times 360°}{\textbf{Key formula:}\quad $\displaystyle\text{\ealkey{sector} \ealkey{angle}} = \dfrac{\text{amount}}{\text{total}} \times 360°$}
\end{mdframed}
\vspace{6pt}

%% ===========================
%%  CHART 1 -- Heavily scaffolded: angles given, 2/5 sectors pre-drawn
%%  Data: total 36 students; angles all multiples of 10
%%  Dog 12 → 120°, Cat 9 → 90°, Fish 8 → 80°, Rabbit 5 → 50°, Other 2 → 20°  (sum=360)
%% ===========================
\begin{center}
\ealpara{چارٹ 1 – پسندیدہ پالتو جانور (36 طلباء کی کلاس)}{\large\bfseries Chart 1 \textendash{} Favourite Pets (class of 36 students)}
\end{center}

\noindent
\begin{minipage}[c]{0.6\textwidth}
  \centering
  % Pre-draw: Dog (120°) and Cat (90°) -- student draws the remaining three
  \partialchart{4.4}{Dog/120/slice1,Cat/90/slice2}
\end{minipage}%
\hfill
\begin{minipage}[c]{0.45\textwidth}
  \setlength{\tabcolsep}{4pt}
  \renewcommand{\arraystretch}{1.8}
  \begin{tabular}{|l|c|c|c|}
  \hline
  \textbf{Pet} & \textbf{No.} & \textbf{\ealkey{Fraction}} & \textbf{\ealkey{Angle}} \\
  \hline
  Dog    & 12 & $\tfrac{1}{3}$ & 120\textdegree \\
  \hline
  Cat    &  9 & $\tfrac{1}{4}$ & 90\textdegree \\
  \hline
  Fish   &  8 & $\tfrac{2}{9}$ & 80\textdegree \\
  \hline
  Rabbit &  5 & $\tfrac{5}{36}$ & 50\textdegree \\
  \hline
  Other  &  2 & $\tfrac{1}{18}$ & 20\textdegree \\
  \hline
  \textbf{Total} & \textbf{36} & \textbf{1} & \textbf{360\textdegree} \\
  \hline
  \end{tabular}
\end{minipage}

\vspace{4pt}
\begin{mdframed}[backgroundcolor=boxorange, linecolor=slice2,
                 linewidth=1pt, roundcorner=5pt, innertopmargin=4pt]
\ealpara{\textbf{شروع کریں:} \textbf{Dog} (120\textdegree) اور \textbf{Cat} (90\textdegree) کے \ealkeytr{قطاع} آپ کے لیے پہلے ہی کھینچے جا چکے ہیں۔ \textbf{Cat} کے \ealkeytr{نصف قطر} سے شروع کرتے ہوئے، اپنے \ealkeytr{زاویہ نما} کا استعمال کرتے ہوئے ترتیب سے \textbf{Fish} (80\textdegree)، \textbf{Rabbit} (50\textdegree) اور \textbf{Other} (20\textdegree) شامل کریں۔ جب آپ ہر \ealkeytr{قطاع} کھینچ لیں تو اس پر لیبل لگائیں۔}{\small\textbf{Getting started:} The \textbf{Dog} (120\textdegree) and \textbf{Cat} (90\textdegree) \ealkey{sectors} have already been drawn for you. Starting from the Cat \ealkey{radius}, use your \ealkey{protractor} to add \textbf{Fish} (80\textdegree), \textbf{Rabbit} (50\textdegree) and \textbf{Other} (20\textdegree) in order. Label each \ealkey{sector} when you have drawn it.}
\end{mdframed}

\vspace{4pt}
\ealpara{سوالات:}{Questions:}
\begin{enumerate}[leftmargin=1.6em, itemsep=3pt]
  \item \ealpara{\ealkeytr{پائی چارٹ} مکمل کریں}{Complete the \ealkey{pie chart}}
  \item \ealpara{کس پالتو جانور کو کلاس کے بالکل $\tfrac{1}{3}$ حصے نے چنا؟}{Which pet was chosen by exactly $\tfrac{1}{3}$ of the class?}
  \item \ealpara{کس ایک پالتو جانور کو بالکل 9 طلباء نے چنا؟}{Which single pet was chosen by exactly 9 students?} \quad \ansbox
\end{enumerate}

\newpage

%% ===========================
%%  CHART 2 -- Scaffolded: angles given, 1 sector pre-drawn
%%  Data: total 36 members; all angles multiples of 10
%%  Football 14→140°, Swimming 9→90°, Tennis 6→60°, Basketball 5→50°, Other 2→20°  (sum=360)
%% ===========================
\begin{center}
\ealpara{چارٹ 2 – اسپورٹس کلب ممبرشپس (36 ممبران)}{\large\bfseries Chart 2 \textendash{} Sports Club Memberships (36 members)}
\end{center}

\noindent
\begin{minipage}[c]{0.52\textwidth}
  \centering
  % Pre-draw: Football (140°) only -- student draws the remaining four
  \partialchart{4.4}{Football/140/slice1}
\end{minipage}%
\hfill
\begin{minipage}[c]{0.45\textwidth}
  \setlength{\tabcolsep}{4pt}
  \renewcommand{\arraystretch}{1.8}
  \begin{tabular}{|l|c|c|c|}
  \hline
  \textbf{Sport} & \textbf{No.} & \textbf{\ealkey{Fraction}} & \textbf{\ealkey{Angle}} \\
  \hline
  Football   & 14 & $\tfrac{7}{18}$ & 140\textdegree \\
  \hline
  Swimming   &  9 & $\tfrac{1}{4}$ & 90\textdegree \\
  \hline
  Tennis     &  6 & $\tfrac{1}{6}$ & 60\textdegree \\
  \hline
  Basketball &  5 & $\tfrac{5}{36}$ & 50\textdegree \\
  \hline
  Other      &  2 & $\tfrac{1}{18}$ & 20\textdegree \\
  \hline
  \textbf{Total} & \textbf{36} & \textbf{1} & \textbf{360\textdegree} \\
  \hline
  \end{tabular}
\end{minipage}

\vspace{4pt}

\vspace{4pt}
\ealpara{سوالات:}{Questions:}
\begin{enumerate}[leftmargin=1.6em, itemsep=3pt]
  \item \ealpara{\ealkeytr{پائی چارٹ} مکمل کریں۔}{Complete the \ealkey{pie chart}.}
  \item \ealpara{کس کھیل کا \ealkeytr{قطاع} \ealkeytr{زاویہ} بالکل 90\textdegree ہے؟}{Which sport has a \ealkey{sector} \ealkey{angle} of exactly 90\textdegree?} \quad \ansbox
  \item \ealpara{کتنے ممبر \textbf{Tennis} کھیلتے ہیں؟ \ealkeytr{کسر} میں لکھیں۔}{What \ealkey{fraction} of members play \textbf{Tennis}?} \quad $\dfrac{\phantom{ABC}}{\phantom{ABC}}$
  \item \ealpara{\textbf{تین سب سے چھوٹے} کھیلوں کے لیے کل کتنے \ealkeytr{ڈگری} لیے جاتے ہیں؟}{How many \ealkey{degrees} are taken up by the \textbf{three smallest} sports combined?}
        \vspace{0.6cm}\\
        \hspace*{1.6em}Answer: \ansbox\textdegree
\end{enumerate}

\newpage

%% ===========================
%%  CHART 3 -- Less scaffolded: angles given in table, blank circle
%% ===========================
\begin{center}
\ealpara{چارٹ 3 – طلباء اپنی جیب خرچ کا زیادہ تر حصہ کیسے خرچ کرتے ہیں (36 طلباء)}{\large\bfseries Chart 3 \textendash{} How Students Spend Most of Their Pocket Money (36 students)}
\end{center}

\noindent
\begin{minipage}[c]{0.52\textwidth}
  \centering
  \emptycircle{4.4}
\end{minipage}%
\hfill
\begin{minipage}[c]{0.45\textwidth}
  \setlength{\tabcolsep}{4pt}
  \renewcommand{\arraystretch}{1.8}
  \begin{tabular}{|l|c|c|c|}
  \hline
  \textbf{Category} & \textbf{No.} & \textbf{\ealkey{Fraction}} & \textbf{\ealkey{Angle}} \\
  \hline
  Food \& Drink & 12 & $\tfrac{1}{3}$ & 120\textdegree \\
  \hline
  Games         &  9 & $\tfrac{1}{4}$ & 90\textdegree \\
  \hline
  Clothes       &  6 & $\tfrac{1}{6}$ & 60\textdegree \\
  \hline
  Savings       &  6 & $\tfrac{1}{6}$ & 60\textdegree \\
  \hline
  Other         &  3 & $\tfrac{1}{12}$ & 30\textdegree \\
  \hline
  \textbf{Total} & \textbf{36} & \textbf{1} & \textbf{360\textdegree} \\
  \hline
  \end{tabular}
\end{minipage}


\vspace{4pt}
\ealpara{سوالات:}{Questions:}
\begin{enumerate}[leftmargin=1.6em, itemsep=3pt]
   \item \ealpara{\ealkeytr{پائی چارٹ} مکمل کریں۔}{Complete the \ealkey{pie chart}.}
  \item \ealpara{دو زمرے ایک جیسے \ealkeytr{قطاع} \ealkeytr{زاویہ} کے حامل ہیں۔ وہ کون سے دو ہیں؟}{Two categories have the same \ealkey{sector} \ealkey{angle}. Which two are they?} \quad \ansbox
\end{enumerate}

\newpage

%% ===========================
%%  CHART 4 -- Partially scaffolded table: No. given, students calculate fraction AND angle
%% ===========================
\begin{center}
\ealpara{چارٹ 4 – کار پارک میں کاروں کے رنگ (60 کاریں)}{\large\bfseries Chart 4 \textendash{} Colours of Cars in a Car Park (60 cars)}
\end{center}

\noindent
\begin{minipage}[c]{0.52\textwidth}
  \centering
  \emptycircle{4.4}
\end{minipage}%
\hfill
\begin{minipage}[c]{0.45\textwidth}
  \setlength{\tabcolsep}{4pt}
  \renewcommand{\arraystretch}{2.0}
  \begin{tabular}{|l|c|c|c|}
  \hline
  \textbf{Colour} & \textbf{No.} & \textbf{\ealkey{Fraction}} & \textbf{\ealkey{Angle}} \\
  \hline
  White  & 20 & & \\
  \hline
  Black  & 15 & & \\
  \hline
  Silver & 12 & & \\
  \hline
  Red    &  9 & & \\
  \hline
  Other  &  4 & & \\
  \hline
  \textbf{Total} & \textbf{60} & \textbf{1} & \textbf{360\textdegree} \\
  \hline
  \end{tabular}
\end{minipage}

\vspace{6pt}
\ealpara{ہدایات:}{Instructions:}
\begin{enumerate}[leftmargin=1.6em, itemsep=2pt]
  \item \ealpara{\ealkeytr{کسر} کا کالم مکمل کریں۔ ہر \ealkeytr{کسر} کو اس کی \ealkeytr{سادہ ترین صورت} میں لکھیں۔}{Complete the \ealkey{Fraction} column. Write each \ealkey{fraction} in its \ealkey{simplest form}.}
  \item \ealpara{\ealkeytr{زاویہ} کا کالم مکمل کرنے کے لیے $\text{angle} = \tfrac{\text{amount}}{\text{total}} \times 360°$ استعمال کریں۔}{Use $\text{\ealkey{angle}} = \tfrac{\text{amount}}{\text{total}} \times 360°$ to complete the \ealkey{Angle} column.}
  \item \ealpara{چیک کریں کہ آپ کے \ealkeytr{زاویہ} کا \ealkeytr{حاصل جمع} 360\textdegree ہے، پھر \ealkeytr{پائی چارٹ} کھینچیں اور لیبل لگائیں۔}{Check your \ealkey{angles} \ealkey{sum} to 360\textdegree, then draw and label the \ealkey{pie chart}.}
\end{enumerate}

\vspace{4pt}
\ealpara{سوالات:}{Questions:}
\begin{enumerate}[leftmargin=1.6em, itemsep=3pt]
\item \ealpara{جدول مکمل کریں۔}{Complete the table.}
\item \ealpara{\ealkeytr{پائی چارٹ} مکمل کریں۔}{Complete the \ealkey{pie chart}.}
  \item \ealpara{\textbf{White} کاروں کی نمائندگی کون سا \ealkeytr{زاویہ} کرتا ہے؟ اپنا کام دکھائیں۔}{What \ealkey{angle} represents \textbf{White} cars? Show your working.}
        \vspace{0.9cm}\\
        \hspace*{1.6em}Answer: \ansbox\textdegree
  \item \ealpara{\textbf{Black} کاروں کا کون سا \ealkeytr{کسر} ہے؟}{What \ealkey{fraction} of cars are \textbf{Black}?} \quad $\dfrac{\phantom{ABC}}{\phantom{ABC}}$
  \item \ealpara{چیک کریں: کیا آپ کے \ealkeytr{زاویہ} کا مجموعہ 360\textdegree ہے؟ اپنا کل لکھیں۔}{Check: do your \ealkey{angles} add up to 360\textdegree? Write your total.} \quad \ansbox\textdegree
\end{enumerate}

\newpage

%% ===========================
%%  CHART 5 -- Least scaffolded: only raw data given
%% ===========================
\begin{center}
\ealpara{چارٹ 5 – گھر میں بولی جانے والی زبانیں (90 طلباء کا سروے)}{\large\bfseries Chart 5 \textendash{} Languages Spoken at Home (survey of 90 students)}
\end{center}

\noindent
\begin{minipage}[c]{0.52\textwidth}
  \centering
  \emptycircle{4.4}
\end{minipage}%
\hfill
\begin{minipage}[c]{0.45\textwidth}
  \setlength{\tabcolsep}{4pt}
  \renewcommand{\arraystretch}{2.0}
  \begin{tabular}{|l|c|c|c|}
  \hline
  \textbf{Language} & \textbf{No.} & \textbf{\ealkey{Fraction}} & \textbf{\ealkey{Angle}} \\
  \hline
  English  & 45 & & \\
  \hline
  Urdu     & 18 & & \\
  \hline
  Polish   &  9 & & \\
  \hline
  Punjabi  &  9 & & \\
  \hline
  Other    &  9 & & \\
  \hline
  \textbf{Total} & \textbf{90} & & \\
  \hline
  \end{tabular}
\end{minipage}

\vspace{6pt}
\ealpara{سوالات:}{Questions:}
\begin{enumerate}[leftmargin=1.6em, itemsep=3pt]
  \item \ealpara{جدول مکمل کریں۔}{Complete the table.}
  \item \ealpara{\ealkeytr{پائی چارٹ} مکمل کریں۔}{Complete the \ealkey{pie chart}.}
  \item \ealpara{تین زبانوں کے بولنے والوں کی تعداد ایک جیسی ہے۔ ان میں سے ہر ایک کے \ealkeytr{قطاع} کا کون سا \ealkeytr{زاویہ} ہے؟}{Three languages have the same number of speakers. What \ealkey{angle} does each of their \ealkey{sectors} have?}
        \vspace{0.6cm}\\
        \hspace*{1.6em}Answer: \ansbox\textdegree
  \item \ealpara{انگریزی بولنے والے سروے کا کون سا \ealkeytr{کسر} بناتے ہیں؟ اسے \ealkeytr{سادہ ترین صورت} میں لکھیں۔}{English speakers make up what \ealkey{fraction} of the survey? Write it in its \ealkey{simplest form}.} \quad $\dfrac{\phantom{ABC}}{\phantom{ABC}}$
  \item \ealpara{\textbf{چیلنج:} اگر 10 مزید طلباء سروے میں شامل ہوں اور وہ سب پولش بولتے ہوں، تو پولش کے لیے نیا \ealkeytr{زاویہ} کیا ہوگا؟ اپنا مکمل کام نیچے دی گئی جگہ میں دکھائیں۔}{\textbf{Challenge:} If 10 more students joined the survey and they all spoke Polish, what would the new \ealkey{angle} for Polish be? Show your full working in the space below.}

  \vspace{10em}

  
  Answer: \ansbox\textdegree
\end{enumerate}

\newpage

% %% ===========================
% %%  ANSWER KEY
% %% ===========================
% \begin{mdframed}[backgroundcolor=lightgrey, linecolor=gray,
%                  linewidth=1pt, roundcorner=4pt, innertopmargin=5pt]
% \textbf{\small Answers (Teacher Copy):}\\[3pt]
% {\small
% \textbf{Chart 1 -- Pets (36 students):}
% Angles: Dog 120\textdegree, Cat 90\textdegree, Fish 80\textdegree, Rabbit 50\textdegree, Other 20\textdegree.
% Q1: Dog ($\tfrac{1}{3}$).
% Q2: $360-120-90=150$\textdegree.
% Q3: Cat (9 students).\\[3pt]
% \textbf{Chart 2 -- Sports (36 members):}
% Angles: Football 140\textdegree, Swimming 90\textdegree, Tennis 60\textdegree, Basketball 50\textdegree, Other 20\textdegree.
% Q1: Swimming (90\textdegree).
% Q2: $\tfrac{1}{6}$.
% Q3: Tennis + Basketball + Other = $60+50+20=130$\textdegree.
% Q4: Yes (or student explains checking sum = 360\textdegree{} and redrawing if not).\\[3pt]
% \textbf{Chart 3 -- Pocket Money:}
% Q1: Clothes \& Savings (both 60\textdegree).
% Q2: Clothes (or Savings) -- both are $\tfrac{1}{6}$.
% Q3: $\tfrac{1}{3}\times48=16$ students.\\[3pt]
% \textbf{Chart 4 -- Cars:}
% White: $\tfrac{20}{60}=\tfrac{1}{3}$, $120$\textdegree.
% Black: $\tfrac{15}{60}=\tfrac{1}{4}$, $90$\textdegree.
% Silver: $\tfrac{12}{60}=\tfrac{1}{5}$, $72$\textdegree.
% Red: $\tfrac{9}{60}=\tfrac{3}{20}$, $54$\textdegree.
% Other: $\tfrac{4}{60}=\tfrac{1}{15}$, $24$\textdegree.
% Sum: $120+90+72+54+24=360$\textdegree\checkmark.
% Q1: 120\textdegree.
% Q2: $\tfrac{1}{4}$.
% Q3: 360\textdegree.\\[3pt]
% \textbf{Chart 5 -- Languages:}
% English: $\tfrac{1}{2}$, 180\textdegree.\;
% Urdu: $\tfrac{1}{5}$, 72\textdegree.\;
% Polish: $\tfrac{1}{10}$, 36\textdegree.\;
% Punjabi: $\tfrac{1}{10}$, 36\textdegree.\;
% Other: $\tfrac{1}{10}$, 36\textdegree.
% Q1: 36\textdegree{} each.
% Q2: $\tfrac{1}{2}$.
% Challenge: New total = 100; Polish = 19; angle = $\tfrac{19}{100}\times360=68.4$\textdegree{} (accept 68\textdegree{} or 69\textdegree).
% }
% \end{mdframed}

\end{document}